% answers.tex —— 1991 年·全国卷·理科 · 参考答案与详细解析片段
% 由 tools/generate_projects.py 依来源总账自 1991/1991_数学真题_全国理_试题.md + 1991/1991_数学真题_全国理_解析正文.md + 1991/1991_数学真题_全国理_答案要点.md 生成，勿手改（改源 MD 后重跑）


\begin{gktip}
本稿根据固定来源 PDF 的题面推导。原件没有答案、解析或分值；所有结论均为非官方推导，尚待独立数学复核。
\end{gktip}

\begin{answer}
\gksolhead{1}{}

\textbf{题目}：

已知 $\sin\alpha=\dfrac45$，并且 $\alpha$ 是第二象限的角，那么 $\tan\alpha$ 的值等于\mbox{（　　）}

（A）$-\dfrac43$　　（B）$-\dfrac34$　　（C）$\dfrac34$　　（D）$\dfrac43$

\Needspace{4\baselineskip}
\textbf{解析}

第二象限内 $\cos\alpha=-\sqrt{1-\sin^2\alpha}=-\dfrac35$，故 $\tan\alpha=\dfrac{\sin\alpha}{\cos\alpha}=-\dfrac43$，选 A。

\gkanswer{答案：A}

\end{answer}
\begin{answer}
\gksolhead{2}{}

\textbf{题目}：

焦点在 $(-1,0)$，顶点在 $(1,0)$ 的抛物线方程是\mbox{（　　）}

（A）$y^2=8(x+1)$　　（B）$y^2=-8(x+1)$

（C）$y^2=8(x-1)$　　（D）$y^2=-8(x-1)$

\Needspace{4\baselineskip}
\textbf{解析}

顶点为 $(1,0)$、焦点为 $(-1,0)$，故抛物线开口向左，$p=-2$。由 $y^2=4p(x-1)$，得 $y^2=-8(x-1)$，选 D。

\gkanswer{答案：D}

\end{answer}
\begin{answer}
\gksolhead{3}{}

\textbf{题目}：

函数 $y=\cos^4x-\sin^4x$ 的最小正周期是\mbox{（　　）}

（A）$\dfrac\pi2$　　（B）$\pi$　　（C）$2\pi$　　（D）$4\pi$

\Needspace{4\baselineskip}
\textbf{解析}

利用平方差，$\cos^4x-\sin^4x=(\cos^2x-\sin^2x)(\cos^2x+\sin^2x)=\cos2x$，最小正周期为 $\pi$，选 B。

\gkanswer{答案：B}

\end{answer}
\begin{answer}
\gksolhead{4}{}

\textbf{题目}：

如果把两条异面直线看成“一对”，那么六棱锥的棱所在的 12 条直线中，异面直线共有\mbox{（　　）}

（A）12 对　　（B）24 对　　（C）36 对　　（D）48 对

\Needspace{4\baselineskip}
\textbf{解析}

六条底棱所在直线共面，六条侧棱所在直线交于锥顶，因此这两类内部都没有异面直线对。每一条侧棱直线与两条相邻底棱直线相交；与其余四条底棱直线异面。故异面直线对数为 $6\times4=24$，选 B。

\gkanswer{答案：B}

\end{answer}
\begin{answer}
\gksolhead{5}{}

\textbf{题目}：

函数 $y=\sin\left(2x+\dfrac{5\pi}{2}\right)$ 的图象的一条对称轴的方程是\mbox{（　　）}

（A）$x=-\dfrac\pi2$　　（B）$x=-\dfrac\pi4$　　（C）$x=\dfrac\pi8$　　（D）$x=\dfrac{5\pi}4$

\Needspace{4\baselineskip}
\textbf{解析}

$\sin(2x+5\pi/2)=\cos2x$。函数 $\cos2x$ 的对称轴为 $x=k\pi/2$（$k\in\mathbb Z$），选项中只有 $x=-\pi/2$ 符合，选 A。

\gkanswer{答案：A}

\end{answer}
\begin{answer}
\gksolhead{6}{}

\textbf{题目}：

如果三棱锥 $S-ABC$ 的底面是不等边三角形，侧面与底面所成的二面角都相等，且顶点 $S$ 在底面的射影 $O$ 在 $\triangle ABC$ 内，那么 $O$ 是 $\triangle ABC$ 的\mbox{（　　）}

（A）垂心　　（B）重心　　（C）外心　　（D）内心

\Needspace{4\baselineskip}
\textbf{解析}

设棱锥高为 $h$，从 $S$ 向底面作垂线，垂足为 $O$。对底面任一边，侧面与底面的二面角的正切等于 $h$ 除以 $O$ 到该边所在直线的距离。三个二面角相等且 $h$ 固定，故 $O$ 到三边所在直线的距离相等。又 $O$ 在三角形内部，所以 $O$ 是内心，选 D。

\gkanswer{答案：D}

\end{answer}
\begin{answer}
\gksolhead{7}{}

\textbf{题目}：

已知 $\{a_n\}$ 是等比数列，且 $a_n>0$，$a_2a_4+2a_3a_5+a_4a_6=25$，那么 $a_3+a_5$ 的值等于\mbox{（　　）}

（A）5　　（B）10　　（C）15　　（D）20

\Needspace{4\baselineskip}
\textbf{解析}

\Needspace{5\baselineskip}
设公比为 $q$。由等比数列性质，$a_2a_4=a_3^2$、$a_3a_5=a_3^2q^2$、$a_4a_6=a_3^2q^4$，于是
\nopagebreak[4]
\[
a_2a_4+2a_3a_5+a_4a_6=a_3^2(1+q^2)^2=(a_3+a_5)^2=25
\]
因 $a_n>0$，故 $a_3+a_5=5$，选 A。

\gkanswer{答案：A}

\end{answer}
\begin{answer}
\gksolhead{8}{}

\textbf{题目}：

如果圆锥曲线的极坐标方程为 $\rho=\dfrac{16}{5-3\cos\theta}$，那么它的焦点的极坐标为\mbox{（　　）}

（A）$(0,0)$、$(6,\pi)$　　（B）$(-3,0)$、$(3,0)$

（C）$(0,0)$、$(3,0)$　　（D）$(0,0)$、$(6,0)$

\Needspace{4\baselineskip}
\textbf{解析}

\Needspace{5\baselineskip}
由 $\rho=\dfrac{16}{5-3\cos\theta}$ 及 $x=\rho\cos\theta$，有 $5\rho=16+3x$。平方并利用 $\rho^2=x^2+y^2$，得
\nopagebreak[4]
\[
25(x^2+y^2)=(16+3x)^2,
\]
整理为 $\dfrac{(x-3)^2}{25}+\dfrac{y^2}{16}=1$。焦点为 $(3\pm3,0)$，即 $(0,0)$、$(6,0)$，选 D。

\gkanswer{答案：D}

\end{answer}
\begin{answer}
\gksolhead{9}{}

\textbf{题目}：

从 4 台甲型和 5 台乙型电视机中任意取出 3 台，其中至少要有甲型与乙型电视机各 1 台，则不同的取法共有\mbox{（　　）}

（A）140 种　　（B）84 种　　（C）70 种　　（D）35 种

\Needspace{4\baselineskip}
\textbf{解析}

从全部 9 台中选 3 台，再减去全为甲型或全为乙型的取法：$\binom93-\binom43-\binom53=84-4-10=70$，选 C。

\gkanswer{答案：C}

\end{answer}
\begin{answer}
\gksolhead{10}{}

\textbf{题目}：

如果 $AC<0$ 且 $BC<0$，那么直线 $Ax+By+C=0$ 不通过\mbox{（　　）}

（A）第一象限　　（B）第二象限　　（C）第三象限　　（D）第四象限

\Needspace{4\baselineskip}
\textbf{解析}

若 $C>0$，由 $AC<0$、$BC<0$ 得 $A<0$、$B<0$；当 $x<0,y<0$ 时，$Ax+By+C>0$，不可能为 0。若 $C<0$，将方程同乘 $-1$ 可归结为同一情形。因此直线不通过第三象限，选 C。

\gkanswer{答案：C}

\end{answer}
\begin{answer}
\gksolhead{11}{}

\textbf{题目}：

设甲、乙、丙是三个命题。如果甲是乙的必要条件；丙是乙的充分条件但不是乙的必要条件，那么\mbox{（　　）}

（A）丙是甲的充分条件，但不是甲的必要条件

（B）丙是甲的必要条件，但不是甲的充分条件

（C）丙是甲的充要条件

（D）丙不是甲的充分条件，也不是甲的必要条件

\Needspace{4\baselineskip}
\textbf{解析}

“甲是乙的必要条件”表示 $\text{乙}\Rightarrow\text{甲}$；“丙是乙的充分条件”表示 $\text{丙}\Rightarrow\text{乙}$，故 $\text{丙}\Rightarrow\text{甲}$。又丙不是乙的必要条件，存在乙真而丙假的情形；此时甲真、丙假，说明丙不是甲的必要条件。因此选 A。

\gkanswer{答案：A}

\end{answer}
\begin{answer}
\gksolhead{12}{}

\Needspace{5\baselineskip}
\textbf{题目}：

\nopagebreak[4]
\[
\lim_{n\to\infty}\left[n\left(1-\frac13\right)\left(1-\frac14\right)\left(1-\frac15\right)\cdots\left(1-\frac1{n+2}\right)\right]
\]

的值等于\mbox{（　　）}

（A）0　　（B）1　　（C）2　　（D）3

\Needspace{4\baselineskip}
\textbf{解析}

\Needspace{5\baselineskip}
括号内的乘积可望远镜化：
\nopagebreak[4]
\[
n\prod_{k=3}^{n+2}\left(1-\frac1k\right)
=n\prod_{k=3}^{n+2}\frac{k-1}{k}
=n\cdot\frac2{n+2}
\]
取极限得 $2$，选 C。

\gkanswer{答案：C}

\end{answer}
\begin{answer}
\gksolhead{13}{}

\textbf{题目}：

如果奇函数 $f(x)$ 在区间 $[3,7]$ 上是增函数且最小值为 5，那么 $f(x)$ 在区间 $[-7,-3]$ 上是\mbox{（　　）}

（A）增函数且最小值为 $-5$　　（B）增函数且最大值为 $-5$

（C）减函数且最小值为 $-5$　　（D）减函数且最大值为 $-5$

\Needspace{4\baselineskip}
\textbf{解析}

奇函数满足 $f(-x)=-f(x)$。当 $x$ 从 $-7$ 增加到 $-3$ 时，$-x$ 从 $7$ 减少到 $3$；由 $f$ 在 $[3,7]$ 上递增可知 $f(x)$ 仍递增。并且 $f(-3)=-f(3)=-5$，所以在 $[-7,-3]$ 上最大值为 $-5$，选 B。

\gkanswer{答案：B}

\end{answer}
\begin{answer}
\gksolhead{14}{}

\textbf{题目}：

圆 $x^2+2x+y^2+4y-3=0$ 上到直线 $x+y+1=0$ 的距离为 $\sqrt2$ 的点共有\mbox{（　　）}

（A）1 个　　（B）2 个　　（C）3 个　　（D）4 个

\Needspace{4\baselineskip}
\textbf{解析}

圆方程化为 $(x+1)^2+(y+2)^2=8$，圆心 $(-1,-2)$，半径 $2\sqrt2$。到直线 $x+y+1=0$ 的距离为 $\sqrt2$，等价于 $|x+y+1|=2$，所以这些点位于两条平行线 $x+y-1=0$ 或 $x+y+3=0$ 上。前者到圆心的距离为 $2\sqrt2$，与圆相切于 1 点；后者经过圆心，与圆相交于 2 点。共有 3 点，选 C。

\gkanswer{答案：C}

\end{answer}
\begin{answer}
\gksolhead{15}{}

\textbf{题目}：

设全集为 $\mathbb R$，$f(x)=\sin x$，$g(x)=\cos x$，$M=\{x\mid f(x)\ne0\}$，$N=\{x\mid g(x)\ne0\}$，那么集合 $\{x\mid f(x)g(x)=0\}$ 等于\mbox{（　　）}

（A）$\overline M\cap\overline N$　　（B）$\overline M\cup N$　　（C）$M\cup\overline N$　　（D）$\overline M\cup\overline N$

\Needspace{4\baselineskip}
\textbf{解析}

$f(x)g(x)=0$ 当且仅当 $\sin x=0$ 或 $\cos x=0$。由定义，$\overline M=\{x\mid\sin x=0\}$、$\overline N=\{x\mid\cos x=0\}$，所求集合为 $\overline M\cup\overline N$，选 D。

\gkanswer{答案：D}

\end{answer}
\begin{answer}
\gksolhead{16}{}

\textbf{题目}：

$\arctan\dfrac13+\arctan\dfrac12$ 的值是\underline{\hspace{6em}}。

\Needspace{4\baselineskip}
\textbf{解析}

\Needspace{5\baselineskip}
因为 $0<\dfrac13<1$、$0<\dfrac12<1$，且反正切函数严格递增、$\arctan1=\dfrac\pi4$，所以
\nopagebreak[4]
\[
0<\arctan\dfrac13<\dfrac\pi4,\qquad
0<\arctan\dfrac12<\dfrac\pi4.
\]
\Needspace{5\baselineskip}
两角之和因此位于 $(0,\dfrac\pi2)$ 内。由正切和角公式，
\nopagebreak[4]
\[
\tan\left(\arctan\dfrac13+\arctan\dfrac12\right)
=\dfrac{\dfrac13+\dfrac12}{1-\dfrac16}=1
\]
故和角为 $\dfrac\pi4$。

\gkanswer{答案：$\dfrac\pi4$}

\end{answer}
\begin{answer}
\gksolhead{17}{}

\textbf{题目}：

不等式 $6^{x^2+x-2}<1$ 的解集是\underline{\hspace{8em}}。

\Needspace{4\baselineskip}
\textbf{解析}

因底数 $6>1$，不等式等价于 $x^2+x-2<0$，即 $(x+2)(x-1)<0$，解集为 $(-2,1)$。

\gkanswer{答案：$(-2,1)$}

\end{answer}
\begin{answer}
\gksolhead{18}{}

\textbf{题目}：

已知正三棱台上底面边长为 2，下底面边长为 4，且侧棱与底面所成的角是 $45^\circ$，那么这个正三棱台的体积等于\underline{\hspace{6em}}。

\Needspace{4\baselineskip}
\textbf{解析}

\Needspace{5\baselineskip}
正三棱台两底面的外接圆半径之差为 $\dfrac4{\sqrt3}-\dfrac2{\sqrt3}=\dfrac2{\sqrt3}$，这是侧棱在底面上的射影长度。侧棱与底面所成角为 $45^\circ$，故台高 $h=2/\sqrt3$。两底面积分别为 $\sqrt3$、$4\sqrt3$，所以
\nopagebreak[4]
\[
V=\frac h3\left(S_1+S_2+\sqrt{S_1S_2}\right)
=\frac{2/\sqrt3}{3}(\sqrt3+4\sqrt3+2\sqrt3)
=\frac{14}3
\]

\gkanswer{答案：$\dfrac{14}3$}

\end{answer}
\begin{answer}
\gksolhead{19}{}

\textbf{题目}：

$(ax+1)^7$ 的展开式中，$x^3$ 的系数是 $x^2$ 的系数与 $x^4$ 的系数的等差中项。若实数 $a>1$，那么 $a=$\underline{\hspace{6em}}。

\Needspace{4\baselineskip}
\textbf{解析}

\Needspace{5\baselineskip}
三个系数依次为 $35a^3$、$21a^2$、$35a^4$。按等差中项条件，
\nopagebreak[4]
\[
2\cdot35a^3=21a^2+35a^4
\]
因 $a>1$，可除以 $7a^2$，得 $5a^2-10a+3=0$，故 $a=1\pm\dfrac{\sqrt{10}}5$。由 $a>1$ 取 $1+\dfrac{\sqrt{10}}5$。

\gkanswer{答案：$1+\dfrac{\sqrt{10}}5$}

\end{answer}
\begin{answer}
\gksolhead{20}{}

\textbf{题目}：

在球面上有四个点 $P$、$A$、$B$、$C$，如果 $PA$、$PB$、$PC$ 两两互相垂直，且 $PA=PB=PC=a$，那么这个球面的面积是\underline{\hspace{6em}}。

\Needspace{4\baselineskip}
\textbf{解析}

以 $P$ 为原点、$PA$、$PB$、$PC$ 的方向为三条坐标轴，则 $A,B,C$ 的坐标可取 $(a,0,0)$、$(0,a,0)$、$(0,0,a)$。外接球心为 $(a/2,a/2,a/2)$，半径平方为 $3a^2/4$，球面面积为 $4\pi\cdot3a^2/4=3\pi a^2$。

\gkanswer{答案：$3\pi a^2$}

\end{answer}
\begin{answer}
\gksolhead{21}{}

\textbf{题目}：

求函数 $y=\sin^2x+2\sin x\cos x+3\cos^2x$ 的最小值，并写出使函数 $y$ 取最小值的 $x$ 的集合。

\Needspace{4\baselineskip}
\textbf{解析}

\Needspace{5\baselineskip}
二次型
\nopagebreak[4]
\[
y=\begin{pmatrix}\sin x&\cos x\end{pmatrix}
\begin{pmatrix}1&1\\1&3\end{pmatrix}
\begin{pmatrix}\sin x\\\cos x\end{pmatrix}
\]
的向量长度为 1。矩阵特征值为 $2\pm\sqrt2$，故最小值为 $2-\sqrt2$。对应特征向量满足 $(\sqrt2-1)\sin x+\cos x=0$，所以 $\tan x=-(\sqrt2+1)$，即 $x=-\dfrac{3\pi}8+k\pi$，$k\in\mathbb Z$。

\gkanswer{答案：最小值为 $2-\sqrt2$；$x=-\dfrac{3\pi}8+k\pi$，$k\in\mathbb Z$。}

\end{answer}
\begin{answer}
\gksolhead{22}{}

\textbf{题目}：

已知复数 $z=1+i$，求复数 $\dfrac{z^2-3z+6}{z+1}$ 的模和辐角的主值。

\Needspace{4\baselineskip}
\textbf{解析}

\Needspace{5\baselineskip}
代入 $z=1+i$，有 $z^2-3z+6=3-i$、$z+1=2+i$，因此
\nopagebreak[4]
\[
\frac{z^2-3z+6}{z+1}
=\frac{3-i}{2+i}=1-i
\]
其模为 $\sqrt2$，辐角的主值为 $-\pi/4$。

\gkanswer{答案：模为 $\sqrt2$，辐角的主值为 $-\dfrac\pi4$。}

\end{answer}
\begin{answer}
\gksolhead{23}{}

\textbf{题目}：

已知 $ABCD$ 是边长为 4 的正方形，$E$、$F$ 分别是 $AB$、$AD$ 的中点，$GC$ 垂直于 $ABCD$ 所在的平面，且 $GC=2$。求点 $B$ 到平面 $EFG$ 的距离。

\Needspace{4\baselineskip}
\textbf{解析}

\Needspace{5\baselineskip}
取坐标 $A=(0,0,0)$、$B=(4,0,0)$、$D=(0,4,0)$、$C=(4,4,0)$，则 $E=(2,0,0)$、$F=(0,2,0)$、$G=(4,4,2)$。平面 $EFG$ 的方程为 $x+y-3z-2=0$。点 $B$ 到该平面的距离为
\nopagebreak[4]
\[
\frac{|4+0-0-2|}{\sqrt{1^2+1^2+(-3)^2}}=\frac2{\sqrt{11}}
\]

\gkanswer{答案：$\dfrac2{\sqrt{11}}$。}

\end{answer}
\begin{answer}
\gksolhead{24}{}

\textbf{题目}：

根据函数单调性的定义，证明函数 $f(x)=-x^3+1$ 在 $(-\infty,+\infty)$ 上是减函数。

\Needspace{4\baselineskip}
\textbf{解析}

\Needspace{5\baselineskip}
任取 $x_1<x_2$，则
\nopagebreak[4]
\[
f(x_1)-f(x_2)=x_2^3-x_1^3
=(x_2-x_1)(x_1^2+x_1x_2+x_2^2)>0
\]
因此 $f(x_1)>f(x_2)$，按单调性定义，$f(x)=-x^3+1$ 在 $\mathbb R$ 上严格递减。

\gkanswer{答案：$f(x)=-x^3+1$ 在 $\mathbb R$ 上严格递减。}

\end{answer}
\begin{answer}
\gksolhead{25}{}

\textbf{题目}：

\Needspace{5\baselineskip}
已知 $n$ 为自然数，实数 $a>1$，解关于 $x$ 的不等式
\nopagebreak[4]
\[
\log_a x-4\log_{a^2}x+12\log_{a^3}x+\cdots+n(-2)^{n-1}\log_{a^n}x
>
\frac{1-(-2)^n}{3}\log_a(x^2-a)
\]

\Needspace{4\baselineskip}
\textbf{解析}

\Needspace{5\baselineskip}
原式中的对数有定义要求 $n\ge1$；若自然数约定包含 $n=0$，则此时 $a^n=1$，以 1 为底的对数无定义，不能套用偶数情形。对 $n\ge1$，定义域要求 $x>0$ 且 $x^2-a>0$，合并为 $x>\sqrt a$。利用 $\log_{a^k}x=\dfrac1k\log_a x$，左侧为
\nopagebreak[4]
\[
\sum_{k=1}^n k(-2)^{k-1}\log_{a^k}x
=\sum_{k=1}^n(-2)^{k-1}\log_a x
=\frac{1-(-2)^n}{3}\log_a x
\]
当 $n$ 为奇数时系数为正，因 $a>1$，不等式等价于 $x>x^2-a$，结合定义域得 $\sqrt a<x<\dfrac{1+\sqrt{1+4a}}2$。当 $n$ 为偶数时系数为负，不等号反向，等价于 $x<x^2-a$，结合定义域得 $x>\dfrac{1+\sqrt{1+4a}}2$。

\gkanswer{答案：原式要求 $n\ge1$，定义域为 $x>\sqrt a$。若 $n$ 为奇数，则 $\sqrt a<x<\dfrac{1+\sqrt{1+4a}}2$；若 $n$ 为偶数，则 $x>\dfrac{1+\sqrt{1+4a}}2$。}

\end{answer}
\begin{answer}
\gksolhead{26}{}

\textbf{题目}：

双曲线的中心在坐标原点 $O$，焦点在 $x$ 轴上，过双曲线右焦点且斜率为 $\sqrt{\dfrac35}$ 的直线交双曲线于 $P$、$Q$ 两点。若 $OP\perp OQ$，$|PQ|=4$，求双曲线的方程。

\Needspace{4\baselineskip}
\textbf{解析}

\Needspace{5\baselineskip}
设双曲线为 $\dfrac{x^2}{a^2}-\dfrac{y^2}{b^2}=1$，令 $c^2=a^2+b^2$、$u=a^2/c^2$，斜率平方为 $m^2=3/5$。过右焦点的直线为 $y=m(x-c)$。将其代入双曲线，设交点横坐标为 $x_1,x_2$，由韦达定理并用 $OP\perp OQ$（即 $x_1x_2+y_1y_2=0$）化简，得
\nopagebreak[4]
\[
8u^2-14u+3=0
\]
\Needspace{5\baselineskip}
根为 $u=\dfrac14$ 或 $u=\dfrac32$；由于 $a^2<c^2$，只能取 $u=\dfrac14$，故 $a^2=c^2/4$、$b^2=3c^2/4$。此时交点横坐标满足 $16x^2+8cx-9c^2=0$，故 $x_1+x_2=-c/2$、$x_1x_2=-9c^2/16$，从而
\nopagebreak[4]
\[
(x_1-x_2)^2=(x_1+x_2)^2-4x_1x_2=\frac52c^2
\]
\Needspace{5\baselineskip}
又 $|PQ|^2=(1+m^2)(x_1-x_2)^2=4^2$，所以 $\dfrac85\cdot\dfrac52c^2=16$，得到 $c^2=4$，从而 $a^2=1$、$b^2=3$。所求双曲线为
\nopagebreak[4]
\[
x^2-\frac{y^2}{3}=1
\]

\gkanswer{答案：$x^2-\dfrac{y^2}3=1$。}

\end{answer}